\par\noindent\begin{minipage}{\linewidth}
\begin{DefinitionBox}{Exponentialfunktion}{real_functions_exponential}
  Für eine Basis $a>0$ heißt
  \[
    f(x)=a^x,\qquad x\in\R,
  \]
  eine \emph{Exponentialfunktion} \EnglishTerm{exponential function}.
  Es gilt $a^0=1$ und $a^{x+h}=a^x a^h$: Ein fester Eingabeschritt $h$
  multipliziert den Funktionswert mit dem festen Faktor $a^h$.
\end{DefinitionBox}
\begin{ExampleBox}[before skip=0.3\baselineskip]{Gleicher Zuwachs oder gleicher Faktor?}{real_functions_linear_exponential}
  Beide Funktionen starten bei $1$ und haben bei $x=1$ den Wert $2$.
  Ergänzen Sie die weiteren Werte.
  \[
    \begin{array}{c|rrrrr}
      x&0&1&2&3&4\\ \hline
      1+x&1&2&\Cloze{1cm}{3}&\Cloze{1cm}{4}&\Cloze{1cm}{5}\\[0.6em]
      2^x&1&2&\Cloze{1cm}{4}&\Cloze{1cm}{8}&\Cloze{1cm}{16}
    \end{array}
  \]
  Bei $1+x$ wird je Schritt $1$ addiert, bei $2^x$ mit $2$ multipliziert.
\end{ExampleBox}
\end{minipage}

\begin{center}
  \CachedTikzFigure{content/resources/chapters/real_functions/tikz/functions_exponentials.tex}
\end{center}

\clearpage
\subsection{Monotonie: Wie sind Eingaben und Ausgaben geordnet?}
\label{subsec:real_functions_monotonicity}
$2^x$ wächst stets, $(1/2)^x$ fällt stets. Dieses Verhalten lässt
sich ohne Kenntnis einer Ableitung durch den Vergleich zweier Stellen
beschreiben.
\par\noindent\begin{minipage}{\linewidth}
\begin{FunctionPropertyBox}{Monotonie}{real_functions_monotonicity}
  Eine Funktion $f$ heißt auf einem Intervall $I$
  \emph{monoton steigend} \EnglishTerm{monotonically increasing},
  wenn für alle $x_1,x_2\in I$ gilt:
  \[
    x_1<x_2\quad\Longrightarrow\quad f(x_1)\le f(x_2).
  \]
  Mit $f(x_1)<f(x_2)$ heißt sie \emph{streng monoton steigend}
  \EnglishTerm{strictly increasing}.
  Ersetzt man $\le$ durch $\ge$ beziehungsweise $<$ durch $>$,
  erhält man \emph{monoton fallend} \EnglishTerm{monotonically decreasing}
  beziehungsweise \emph{streng monoton fallend} \EnglishTerm{strictly decreasing}.
  Sie heißt \emph{konstant} \EnglishTerm{constant}, wenn
  $f(x_1)=f(x_2)$ für alle $x_1,x_2\in I$ gilt.
\end{FunctionPropertyBox}
\begin{ExampleBox}[before skip=0.3\baselineskip]{Das Intervall gehört zur Aussage}{real_functions_monotonicity_intervals}
  Betrachten Sie erneut $f(x)=x^2$. Auf welchem Intervall ist die
  Funktion streng monoton fallend, auf welchem streng monoton steigend?
  \[
    \begin{array}{c@{\qquad}c}
      (-\infty,0]&[0,\infty)\\
      \Cloze{5cm}{\text{streng monoton fallend}}&
      \Cloze{5cm}{\text{streng monoton steigend}}
    \end{array}
  \]
  Auf ganz $\R$ ist sie weder monoton steigend noch monoton fallend.
\end{ExampleBox}
\end{minipage}

% Handschrift: PDF-Seite 31 (steigend, konstant und fallend nebeneinander).
\begin{center}
  \CachedTikzFigure{content/resources/chapters/real_functions/tikz/functions_monotonicity_comparison.tex}
\end{center}
Eine konstante Funktion ist sowohl monoton steigend als auch monoton
fallend, aber auf einem Intervall mit mehr als einer Stelle niemals streng monoton.

\begin{TheoremBox}{Eigenschaften der Exponentialfunktion}{real_functions_exponential_properties}
  Für $a>0$, $a\ne1$ ist $a^x$ auf $\R$ stetig, stets positiv und
  hat die \textcolor{green!55!black}{Wertemenge} $(0,\infty)$.
  \begin{center}
    \renewcommand{\arraystretch}{1.4}
    \begin{tabular}{@{}lcc@{}}
      \toprule
       & $a>1$ & $0<a<1$\\ \midrule
      Monotonie & streng steigend & streng fallend\\
      $x\to+\infty$ & $a^x\to+\infty$ & $a^x\to0$\\
      $x\to-\infty$ & $a^x\to0$ & $a^x\to+\infty$\\
      \bottomrule
    \end{tabular}
  \end{center}
  Die Gerade $y=0$ ist in der jeweiligen Richtung eine waagerechte
  Asymptote. Die Funktion ist nach unten, aber nicht nach oben beschränkt.
  Es gibt keine Nullstelle; der $y$-Achsenabschnitt ist $f(0)=a^0=1$.
  Für $a=1$ ist $f(x)=1$ konstant mit \textcolor{green!55!black}{Wertemenge} $\{1\}$.
\end{TheoremBox}
Für $x_2>x_1$ gilt $a^{x_2}/a^{x_1}=a^{x_2-x_1}$.
Dieser Faktor ist bei $a>1$ größer als $1$, bei $0<a<1$ kleiner als $1$;
damit ist die jeweilige strenge Monotonie erklärt.
